\chapter{مرور کارهای مشابه}
\label{chap:related}

برای جایگاه‌یابی چارچوب، خانواده‌های اصلی ابزارهای تست خودکار رابط کاربری مرور می‌شوند و تفاوت رویکرد این پایان‌نامه با آن‌ها بیان می‌گردد.

\section{ابزارهای تست خودکار وب}
\lr{Selenium}
\cite{selenium}
کهن‌ترین و گسترده‌ترین استاندارد خودکارسازی مرورگر است که بر پایه‌ی پروتکل
\lr{WebDriver}
کار می‌کند.
\lr{Cypress}
\cite{cypress}
و
\lr{Playwright}
\cite{playwright}
نسل جدیدترند که با اجرای تست در همان فرایند یا با کنترل مستقیم مرورگر، سرعت و پایداری بهتری ارائه می‌دهند. این ابزارها برای برنامه‌های وب ساخته شده‌اند و به درخت عناصر مرورگر دسترسی دارند؛ برای برنامه‌های دسکتاپ مبتنی بر
\lr{Qt}
کاربرد مستقیمی ندارند. با این حال ایده‌های آن‌ها \pdash{} پروتکل استاندارد درایور، جداسازی لایه‌ی مکان‌یابی از تست‌ها و گزارش‌گیری غنی \pdash{} الهام‌بخش طراحی چارچوب حاضر بوده است.

\section{ابزارهای تجاری تست دسکتاپ}
برای دسکتاپ، ابزارهای تجاری متعددی وجود دارند؛ از جمله
\lr{Ranorex}
و
\lr{TestComplete}
که با بازشناسی اشیاء روی انواع فناوری‌های رابط کاربری کار می‌کنند. برای
\lr{Qt}
به‌طور خاص، ابزار
\lr{Squish}
\cite{squish}
شناخته‌شده‌ترین گزینه است که از طریق بارگذاری کتابخانه در خود برنامه، به اشیاء داخلی
\lr{Qt}
دسترسی عمیق دارد. نقطه ضعف این خانواده هزینه‌ی مجوز، بسته‌بودن متن و وابستگی به فروشنده است؛ همچنین دسترسی عمیق به اشیاء داخلی، برنامه را از دید کاربر واقعی دور می‌کند و رفتار برنامه را در شرایط واقعی نمایش نمی‌دهد.

\section{ابزارهای متن‌باز دسکتاپ و ویندوز}
در سمت متن‌باز،
\lr{WinAppDriver}
\cite{winappdriver}
از طریق
\lr{\gls{UIA}}
\cite{uia}
با هر برنامه‌ی ویندوزی \pdash{} بدون نیاز به تغییر در خود برنامه \pdash{} کار می‌کند و پروتکل آن همان
\lr{WebDriver}
است؛ اما به تنهایی فاقد لایه‌های معماری است: مکان‌یابی، نهان‌گاه دسته‌ها، مدیریت چرخه‌ی پنجره‌ها و گزارش‌گیری بر عهده‌ی نویسنده‌ی تست می‌ماند.
\lr{Qt Test}
\cite{qtcitest}
نیز چارچوب تست داخلی
\lr{Qt}
است که در سطح واحد و یکپارچگی (و حتی تست رابط کاربری درون‌فرایندی) قوی است، اما برنامه را از بیرون و به شکل کاربر واقعی تست نمی‌کند.

\section{جمع‌بندی و جایگاه چارچوب حاضر}
\begin{table}[H]
\centering
\caption{مقایسه‌ی رویکردهای تست خودکار رابط کاربری}
\label{tab:compare}
\begin{tabularx}{\textwidth}{@{}lXXX@{}}
\hline
\lr{Tool} & متن‌باز & دسترسی & دید تست \\
\hline
\lr{Selenium/Cypress/Playwright} & بله & وب & مرورگر \\
\lr{Squish} & خیر & اشیاء داخلی & درون‌فرایندی \\
\lr{Qt Test} & بله & درون‌فرایندی & سطح واحد \\
\lr{WinAppDriver} تنها & بله & \lr{\gls{UIA}} & بیرونی \\
\textbf{چارچوب حاضر} & بله & \lr{\gls{UIA}} + لایه‌ها & بیرونی، کاربر واقعی \\
\hline
\end{tabularx}
\end{table}

چنان‌که در جدول
\autoref{tab:compare}
دیده می‌شود، خلأ موجود، ترکیب سه چیز است: دسترسی بیرونی و بدون تغییر به برنامه، لایه‌های معماری (اشیاء صفحه، زبان مخصوص دامنه، جداسازی داده و گزارش‌گیری) و متن‌باز بودن کامل. چارچوب حاضر دقیقاً این ترکیب را هدف گرفته است: از
\lr{WinAppDriver}
و
\lr{Appium}
به عنوان زیرساخت استفاده می‌کند و آنچه که این زیرساخت نمی‌دهد \pdash{} پایداری، خوانایی و نگهداشت‌پذیری تست‌ها \pdash{} را با لایه‌های معماری و قراردادهای پایداری می‌سازد.